\documentclass{beamer}
\usepackage[T2A]{fontenc}        
\usepackage[utf8]{inputenc}   
\usepackage[russian]{babel}
\usepackage{graphicx}
\usepackage{amsmath}
\usepackage{amssymb}
\usepackage{multirow} 
\title{Презентация на LaTex}
\author{Гаджиева Сара Фазиль кызы}
\date{\today}
\usetheme{CambridgeUS}

\begin{document}
\begin{frame}
\titlepage
\end{frame}

\section{Формулы}
\begin{frame}{Формулы}
Деление комплексных чисел путём умножения на сопряжённое число.
\begin{equation}
\frac{a+bi}{c+di}=\frac{a+bi}{c+di}\cdot\frac{c-di}{c-di}=\frac{ac+bd}{c^2+d^2}+\left(\frac{bc-ad}{c^2+d^2}\right)i
\end{equation}

Интеграл функции, содержащей квадратный корень.
\begin{equation}
\int \frac{du}{u\sqrt{a^2-u^2}}=-\frac{1}{a}\ln\left(\frac{a+\sqrt{a^2-u^2}}{u}\right)
\end{equation}

Производные тригонометрических функций.
\begin{equation}
(\sin x)'=\cos x, \quad (\cos x)'=-\sin x, \quad (\tan x)'=\frac{1}{\cos^2 x}
\end{equation}

\end{frame}


\section{Таблица}
\begin{frame}{Таблица}
\begin{table}[!h]
\centering
\caption{Таблица производных элементарных функций}
\label{tab:1}
\vspace{15pt}
\begin{tabular}{|c|c|c|}
\hline
\textbf{Вид функции} & \textbf{Функция} & \textbf{Ее производная} \\
\hline
\multirow{3}{*}{Алгебраические} 
& $x^n$ & $nx^{n-1}$ \\ \cline{2-3}
& $\frac{1}{x}$ & $-\frac{1}{x^2}$ \\ \cline{2-3}
& $\sqrt{x}$ & $\frac{1}{2\sqrt{x}}$ \\
\hline
\multirow{3}{*}{Показательные и логарифмические}
& $e^x$ & $e^x$ \\ \cline{2-3}
& $\ln x$ & $\frac{1}{x}$ \\ \cline{2-3}
& $a^x$ & $a^x \ln a$ \\
\hline
\multirow{3}{*}{Тригонометрические}
& $\sin x$ & $\cos x$ \\ \cline{2-3}
& $\cos x$ & $-\sin x$ \\ \cline{2-3}
& $\tan x$ & $\frac{1}{\cos^2 x}$ \\
\hline
\end{tabular}
\end{table}
\end{frame}

\section{Списки}
\begin{frame}{Ненумерованный список}
Свойства определенного интеграла
\begin{itemize}
\item линейность:
$\int_a^b (c_1 f+c_2 g) dx=c_1 \int_a^b f dx+c_2 \int_a^b g dx;$
\item аддитивность:
$\int_a^b f dx=\int_a^c f dx+\int_c^b f dx;$
\item интеграл от константы:
$\int_a^b C dx=C(b-a);$
\item свойство монотонности.
\end{itemize}
\end{frame}

\begin{frame}{Нумерованный список}
Основные этапы исследования функции:
\begin{enumerate}
\item Нахождение области определения функции 
\item Исследование функции на четность/нечетность	
\item Поиск точек пересечения с осями координат
\item Нахождение производной функции и критических точек
\item Построение графика функции
\end{enumerate}
\end{frame}

\section{Список}
\begin{frame}{Список литературы}
\begin{thebibliography}{9}
\bibitem{art1} Mathematical Handbook of Formulas and Tables (Third Edition) 
\bibitem{art2} Standard Mathematical tables and formulae (31st edition) - Daniel Zwillinger
\bibitem{art3} Handbook of Mathematical Functions - Abramowitz \& Stegun
\end{thebibliography}
\end{frame}

\end{document}
